% Version 14 of the author's cas-sc copy: every change against the copy as pasted on 2026-10-27.
% Each block: the earlier lines as comments (% OLD:), then the new lines. Files: main.tex and figs/scheme.tex.

% ===== main.tex, old lines 30-30, new lines 30-30
% OLD: \usepackage[authoryear,longnamesfirst]{natbib}
\usepackage[numbers,sort&compress]{natbib}

% ===== main.tex, old lines 241-241, new lines 241-241
% OLD: and Re \cite{SirianniRe2026} are explicit Godunov-type. They
and Re \cite{SirianniRe2026} are explicit Godunov-type schemes. They

% ===== main.tex, old lines 247-247, new lines 247-247
% OLD: all variables. whose convergence is
all variables, whose convergence is

% ===== main.tex, old lines 254-254, new lines 254-254
% OLD: versions split the flux into a convective part with flow-velocity eigenvalues and a pressure part carrying the acoustic
versions split the flux into a convective part, whose eigenvalues are the flow velocity or zero, and a pressure part carrying the acoustic

% ===== main.tex, old lines 315-318, new lines 315-319
% OLD: The discretisation follows a Zha--Bilgen splitting \cite{ZhaBilgen1993}. The explicit step transports all quantities
% OLD: with the flow velocity by the path-conservative Rusanov scheme \cite{Pares2006,Castro2006} with material-speed
% OLD: dissipation. The implicit step carries the pressure, its work on each phase and an incremental flux of the
% OLD: conservative transported quantities. Its central operator is supplemented by Laplace dissipation of pressure and
The discretisation follows the Toro--V\'azquez-Cend\'on splitting \cite{ToroVazquez2012}, with the mass and internal energy
fluxes evaluated at the new velocity. The explicit step transports all quantities with the flow velocity by the
path-conservative Rusanov scheme \cite{Pares2006,Castro2006} with material-speed dissipation. The implicit step carries
the pressure, its work on each phase and the increments of the mass and internal energy fluxes from the old to the new
velocity. Its central operator is supplemented by Laplace dissipation of pressure and

% ===== main.tex, old lines 490-496, new lines 491-496
% OLD: To handle the different time scales of the model we split the fluxes in the manner of Zha and Bilgen
% OLD: \cite{ZhaBilgen1993}, which allows a different time discretisation for each sub-system according to its stiffness: every
% OLD: quantity is transported by the flow velocity in one sub-system, and the pressure, its work and the relaxation terms form
% OLD: the other. The terms of \eqref{eq:general} are separated into the advective part, which transports every component of
% OLD: $\mathbf U$ with the flow velocity, including the non-conservative advection $u\,\partial_x\alpha_1$ of the volume fraction, and the
% OLD: pressure--relaxation part, which contains the pressure flux $\mathbf F^p$, the pressure work of each phase $\mathbf B^p\partial_xu$ and the relaxation terms.
% OLD: Explicitly,
To handle the different time scales of the model we split the fluxes in the manner of Toro and V\'azquez-Cend\'on
\cite{ToroVazquez2012}, which allows a different time discretisation for each sub-system according to its stiffness. The terms
of \eqref{eq:general} are separated into the advective part, which transports the masses, the momentum and the kinetic
energy with the flow velocity and contains the non-conservative advection $u\,\partial_x\alpha_1$ of the volume fraction, and the
pressure--relaxation part, which contains the pressure flux, the internal energy fluxes of the phases with their pressure work
$\mathbf B^p\partial_xu$, the enthalpy flux of the mixture and the relaxation terms. Explicitly,

% ===== main.tex, old lines 499-499, new lines 499-499
% OLD: \mathbf F^c(\mathbf U)=\begin{bmatrix}0\\ m_1u\\ m_2u\\ \rho u^2\\ \varepsilon_1u\\ \varepsilon_2u\\ Eu\end{bmatrix},\quad
\mathbf F^c(\mathbf U)=\begin{bmatrix}0\\ m_1u\\ m_2u\\ \rho u^2\\ 0\\ 0\\ \tfrac12\rho u^3\end{bmatrix},\quad

% ===== main.tex, old lines 501-501, new lines 501-501
% OLD: \mathbf F^p(\mathbf U)=\begin{bmatrix}0\\0\\0\\ \bar p\\0\\0\\ \bar pu\end{bmatrix},\quad
\mathbf F^p(\mathbf U)=\begin{bmatrix}0\\0\\0\\ \bar p\\ \varepsilon_1u\\ \varepsilon_2u\\ \big(E-\tfrac12\rho u^2+\bar p\big)u\end{bmatrix},\quad

% ===== main.tex, old lines 528-528, new lines 528-528
% OLD: \partial_t\mathbf Q+\mathbf A^c(\mathbf Q)\,\partial_x\mathbf Q=\mathbf 0,\
\partial_t\mathbf Q+\mathbf A^c(\mathbf Q)\,\partial_x\mathbf Q=\mathbf 0,\\

% ===== main.tex, old lines 534-538, new lines 534-537
% OLD: 0&0&0&\rho_1\kappa_1&u&0&0\\
% OLD: 0&0&0&\rho_2\kappa_2&0&u&0\\
% OLD: 0&0&0&K^c&0&0&u
% OLD: \end{bmatrix},\qquad
% OLD: K^c=\frac{\sum_k\alpha_k\rho_k\kappa_k/\Gamma_k}{\sum_k\alpha_k/\Gamma_k},
\ast&\ast&0&\rho_1(\kappa_1-\Gamma_1e_1)&0&0&0\\
\ast&0&\ast&\rho_2(\kappa_2-\Gamma_2e_2)&0&0&0\\
\ast&\ast&\ast&\dfrac{\sum_k\alpha_k\rho_k\kappa_k/\Gamma_k-E+\tfrac12\rho u^2}{\sum_k\alpha_k/\Gamma_k}&0&0&0
\end{bmatrix},

% ===== main.tex, old lines 542-545, new lines 541-546
% OLD: since the volume fraction, the velocity and the specific internal energies are advected, the phases are compressed by
% OLD: the flow at constant specific internal energy, and $E-\tfrac12\rho u^2$, transported conservatively, determines $p$ at the advected
% OLD: volumes and compressed masses, with $\kappa_k$ and $\Gamma_k$ evaluated at $\rho_k$ and the pressure $p$; for stiffened and Noble--Abel stiffened gases $K^c=\sum_k\alpha_k(p+\gamma_k\pi_k)/(\gamma_k-1)\big/\sum_k(\alpha_k-m_kb_k)/(\gamma_k-1)$.
% OLD: It admits the eigenvalues
since the volume fraction and the velocity are advected, the masses are transported at fixed phase internal energies
$\varepsilon_k$, and the kinetic energy is transported with the mass, so that $E-\tfrac12\rho u^2$ is constant and determines $p$ at the
advected volumes and transported masses; $\kappa_k$, $\Gamma_k$ and $e_k$ are evaluated at the phase state $\rho_k$, $p_k$ in the rows of $p_k$
and at $\rho_k$ and the pressure $p$ in the last row. The entries $\ast$ are proportional to $u$. For stiffened and Noble--Abel stiffened gases $\kappa_k-\Gamma_ke_k=\Gamma_k\big(b_k(p_k+\gamma_k\pi_k)/(\gamma_k-1)-q_k\big)$ and
the last row has $\sum_km_k\big(b_k(p+\gamma_k\pi_k)/(\gamma_k-1)-q_k\big)\big/\sum_k(\alpha_k-m_kb_k)/(\gamma_k-1)$; both vanish for ideal and stiffened gases
with $q_k=0$. The last three columns of $\mathbf A^c$ vanish, and the eigenvalues are

% ===== main.tex, old lines 547-547, new lines 548-548
% OLD: \lambda^c_{1,\dots,7}=u .
\lambda^c_{1,\dots,4}=u,\qquad \lambda^c_{5,6,7}=0 .

% ===== main.tex, old lines 550-553, new lines 551-553
% OLD: $\mathbf A^c-u\,\mathbf I$ is nilpotent of rank one, so one eigenvector is missing and the sub-system is weakly hyperbolic, as
% OLD: the pressureless gas dynamics it contains. The characteristic speed is the flow velocity alone. The volume fraction
% OLD: is transported only by this sub-system; in the pressure--relaxation sub-system it changes only through the relaxation
% OLD: terms.
One eigenvector of the eigenvalue $u$ is missing, and the sub-system is weakly hyperbolic, as the pressureless gas dynamics it contains.
The characteristic speeds are the flow velocity and zero. The volume fraction is transported only by this sub-system; in
the pressure--relaxation sub-system it changes only through the relaxation terms.

% ===== main.tex, old lines 562-562, new lines 562-562
% OLD: \partial_t\mathbf Q+\mathbf A^p(\mathbf Q)\,\partial_x\mathbf Q=\tilde{\mathbf S}(\mathbf Q),
\partial_t\mathbf Q+\mathbf A^p(\mathbf Q)\,\partial_x\mathbf Q=\tilde{\mathbf S}(\mathbf Q),\\

% ===== main.tex, old lines 568-570, new lines 568-570
% OLD: 0&0&0&\Gamma_1p_1&0&0&0\\
% OLD: 0&0&0&\Gamma_2p_2&0&0&0\\
% OLD: 0&0&0&\bar p\big/\sum_k(\alpha_k/\Gamma_k)&0&0&0
\ast&\ast&0&\rho_1\Gamma_1h_1&u&0&0\\
\ast&0&\ast&\rho_2\Gamma_2h_2&0&u&0\\
\ast&\ast&\ast&\dfrac{E-\tfrac12\rho u^2+\bar p}{\sum_k\alpha_k/\Gamma_k}&0&0&u

% ===== main.tex, old lines 575-578, new lines 575-580
% OLD: with $\tilde{\mathbf S}=(\partial\mathbf Q/\partial\mathbf U)\,\mathbf S$. The fourth row is $\partial_x\bar p/\rho$, $\bar p=\alpha_1p_1+\alpha_2p_2$; each phase keeps its
% OLD: density and is heated by its pressure work, $\partial_tp_k=-\Gamma_kp_k\,\partial_xu$, and the internal energy $E-\tfrac12\rho u^2$ receives $-\bar p\,\partial_xu$, the flux $\partial_x(\bar pu)$ of the total energy
% OLD: minus the work $u\,\partial_x\bar p$ on the kinetic energy, at fixed volumes and phase densities. Only the velocity and the phase pressures are
% OLD: coupled, so the characteristic polynomial of $\mathbf A^p$ is $\lambda^5(\lambda^2-c_{\rm w}^2)$ and the eigenvalues are
with $\tilde{\mathbf S}=(\partial\mathbf Q/\partial\mathbf U)\,\mathbf S$, the coefficients evaluated as in $\mathbf A^c$ and the entries $\ast$ again proportional to $u$. The fourth row is $\partial_x\bar p/\rho$,
$\bar p=\alpha_1p_1+\alpha_2p_2$; each phase keeps its volume and density, and its pressure changes by the transport of $\varepsilon_k$ and by
the pressure work, with the coefficient $\rho_k\Gamma_kh_k$ of $\partial_xu$, $h_k$ the specific enthalpy; the internal energy $E-\tfrac12\rho u^2$ receives
the enthalpy flux of the mixture minus the work $u\,\partial_x\bar p$ on the kinetic energy. The first three rows vanish and the last
column has only its diagonal entry, so the characteristic polynomial of $\mathbf A^p$ is $\lambda^3(\lambda-u)^2(\lambda^2-u\lambda-c_{\rm h}^2)$ and the
eigenvalues are

% ===== main.tex, old lines 580-580, new lines 582-585
% OLD: \lambda^p_{1,\dots,5}=0,\qquad \lambda^p_{6,7}=\pm c_{\rm w},\qquad c_{\rm w}^2=\frac1\rho\sum_k\Gamma_k\alpha_kp_k=\sum_kY_k\,\frac{\Gamma_kp_k}{\rho_k} .
\begin{gathered}
\lambda^p_{1,2,3}=0,\qquad\lambda^p_{4,5}=u,\qquad \lambda^p_{6,7}=\tfrac12\Big(u\pm\sqrt{u^2+4c_{\rm h}^2}\Big),\\
c_{\rm h}^2=\sum_kY_k\Gamma_kh_k .
\end{gathered}

% ===== main.tex, old lines 583-587, new lines 588-596
% OLD: The remaining part of the sound speed, $\sum_kY_k\kappa_k$, comes from the compression of the phases by the mass fluxes of the
% OLD: advective sub-system; the two parts add up to the frozen sound speed, $c_{\rm w}^2+\sum_kY_k\kappa_k=c_f^2$. The
% OLD: speed $c_{\rm w}$ of the pressure work is small in a liquid at moderate pressure, and $c_{\rm w}^2$ is negative for negative phase
% OLD: pressures, as in cavitation; the sub-system is never solved on its own but together with the incremental advective fluxes
% OLD: described below, so that the implicit step carries the whole sound speed, $c_f$ or, with relaxation, that of the active closure.
For $c_{\rm h}^2>0$, $\lambda^p_6>0>\lambda^p_7$ whatever the flow velocity, as for the Euler equations \cite{ToroVazquez2012}. For ideal and
stiffened gases with $q_k=0$, $\Gamma_kh_k=c_k^2$, and the pressure--relaxation sub-system carries the whole frozen sound speed. In
general $c_{\rm h}^2+\sum_kY_k(\kappa_k-\Gamma_ke_k)=c_f^2$, the second part coming from the compression of the phases by the mass fluxes
of the advective sub-system. Through $h_k$, $c_{\rm h}$ depends on the reference energies of the phases, $q_k$ or $\varepsilon_r$: for the water of
Sec.~\ref{sec:res_cav} in its initial state $c_{\rm h}=430$\,m/s against $c_f=1569$\,m/s. Where $c_{\rm h}^2<0$, as in this liquid water below $275$\,K
at $10^5$\,Pa and in the low-pressure initial state of the nitromethane of Sec.~\ref{sec:res_shock}, $\lambda^p_{6,7}$ are complex for $u^2<-4c_{\rm h}^2$
and the sub-system alone is not hyperbolic. It is never solved on its own but together with the incremental mass fluxes
described below, so that the implicit step carries the whole sound speed, $c_f$ or, with relaxation, that of the active
closure.

% ===== main.tex, old lines 592-598, new lines 601-609
% OLD: and the pressure--relaxation sub-system implicitly. So that the implicit step contains the whole acoustic coupling,
% OLD: the advective fluxes are evaluated with the velocity of the new time level, written as the old velocity $u^n$ plus an
% OLD: increment, $Xu^{n+1}=Xu^n+X(u^{n+1}-u^n)$ for $X\in\{m_1,m_2,\varepsilon_1,\varepsilon_2\}$, the momentum and the kinetic
% OLD: energy being carried with the new velocity by the incremental mass flux $\sum_km_k(u^{n+1}-u^n)$: the first part is the
% OLD: explicit advective flux, and the incremental fluxes, which depend on the unknown velocity, are solved together with the
% OLD: pressure--relaxation sub-system in Sec.~\ref{sec:implicit}. Since only the advective eigenvalues enter the explicit stability condition, the time
% OLD: step depends on the flow velocity alone, independently of the sound speed and of the relaxation rates.
and the pressure--relaxation sub-system implicitly. As Chen et al.\ \cite{Chen2026} evaluate the enthalpy flux as $h^n(\rho u)^{n+1}$,
the internal energy fluxes are evaluated with the internal energies at $t^n$ and the velocity at $t^{n+1}$, $\varepsilon^n_ku^{n+1}$, the internal
energy flux of the mixture as their sum and the pressure fluxes as $\bar p^{\,n+1}$ and $\bar p^{\,n+1}u^{n+1}$. So that the implicit step contains the whole
acoustic coupling, the mass fluxes are evaluated in the same way, $m^n_ku^{n+1}$; the momentum and the kinetic energy are carried
with the new velocity by the incremental mass flux $\sum_km_k(u^{n+1}-u^n)$. These fluxes are split into an old-velocity part and an
increment, $Xu^{n+1}=Xu^n+X(u^{n+1}-u^n)$ for $X\in\{m_1,m_2,\varepsilon_1,\varepsilon_2\}$: the first part is evaluated in the explicit step together with the advective fluxes, so that the masses and the internal energies receive the same
numerical dissipation, and the increments, which depend on the unknown velocity, are solved together with the
pressure--relaxation sub-system in Sec.~\ref{sec:implicit}. Since only the flow velocity enters the explicit step, the time step depends
on the flow velocity alone, independently of the sound speed and of the relaxation rates.

% ===== main.tex, old lines 607-608, new lines 618-619
% OLD: After the notation, Sec.~\ref{sec:advective} presents the explicit update of the
% OLD: advective sub-system and Sec.~\ref{sec:implicit} the implicit update of the pressure--relaxation sub-system with its
After the notation, Sec.~\ref{sec:advective} presents the explicit
transport and Sec.~\ref{sec:implicit} the implicit update of the pressure--relaxation sub-system with its

% ===== main.tex, old lines 624-624, new lines 635-635
% OLD: \subsection{Advective sub-system}
\subsection{Explicit transport}

% ===== main.tex, old lines 626-627, new lines 637-640
% OLD: The advective sub-system \eqref{eq:advective} transports every quantity with the flow velocity. Its $\alpha_1$ row,
% OLD: $u\,\partial_x\alpha_1$, is not the derivative of any flux, and the sub-system is discretised by the path-conservative Rusanov scheme
The explicit step advances the advective sub-system \eqref{eq:advective} together with the old-velocity part of the internal
energy fluxes of Sec.~\ref{sec:split}, $\varepsilon_ku^n$ and $(E-\tfrac12\rho u^2)u^n$, so that every quantity is transported with the flow velocity
by the flux $\mathbf F^u(\mathbf U)=u\,(0,m_1,m_2,\rho u,\varepsilon_1,\varepsilon_2,E)^{\mathsf T}$. The $\alpha_1$ row, $u\,\partial_x\alpha_1$, is not the derivative of
any flux, and the step is discretised by the path-conservative Rusanov scheme

% ===== main.tex, old lines 632-632, new lines 645-645
% OLD: \mathbf D^-_{i+1/2}+\mathbf D^+_{i+1/2}=\mathbf F^c(\mathbf U^R)-\mathbf F^c(\mathbf U^L)+\int_0^1\mathbf B^c(\boldsymbol\Psi)\,\partial_\ell\boldsymbol\Psi\,\mathrm d\ell ,
\mathbf D^-_{i+1/2}+\mathbf D^+_{i+1/2}=\mathbf F^u(\mathbf U^R)-\mathbf F^u(\mathbf U^L)+\int_0^1\mathbf B^c(\boldsymbol\Psi)\,\partial_\ell\boldsymbol\Psi\,\mathrm d\ell ,

% ===== main.tex, old lines 642-642, new lines 655-655
% OLD: the eigenvalue $u$ of the advective sub-system on the two face states, or its maximum over all faces of the same
the larger flow speed of the two face states, or its maximum over all faces of the same

% ===== main.tex, old lines 660-661, new lines 673-685
% OLD: and $G^e$ likewise with $v$. For the row of $\alpha_1$ the path integral is taken along the linear path with the face velocity frozen at the mean
% OLD: $\bar u_{i+1/2}=\tfrac12(u^L+u^R)$, $\int_0^1u\,\partial_\ell\alpha_1\,\mathrm d\ell=\bar u_{i+1/2}\,(\alpha_{1,i+1}-\alpha_{1,i})$, and the fluctuation is split
and $G^e$ likewise with $v$. With the mean face velocity $\bar u_{i+1/2}=\tfrac12(u^L+u^R)$ the speed \eqref{eq:rusanov_speed} is
$a_{i+1/2}=|\bar u_{i+1/2}|+\tfrac12|u^R-u^L|$, and \eqref{eq:rusanov} is identically
\begin{equation}
F^e_{i+1/2}[X]=X^{\uparrow}\,\bar u_{i+1/2}-\tfrac12\big(u^L-u^R\big)^+\big(X^R-X^L\big),
\label{eq:rusanov_upwind}
\end{equation}
with $X^\uparrow=X^L$ for $\bar u_{i+1/2}\ge0$, $X^\uparrow=X^R$ otherwise, and $(\cdot)^+=\max(\cdot,0)$. The first term transports $X$ with the mean face velocity
from its upwind side, the second is a dissipation acting only at compressive faces, $u^L>u^R$. With the maximum $a_\Omega$ of
\eqref{eq:rusanov_speed} over the faces of the same direction, the dissipation coefficient is
$\tfrac12(a_\Omega-|\bar u_{i+1/2}|)-\tfrac14(u^R-u^L)\ge\tfrac12(u^L-u^R)^+$.
For the row of $\alpha_1$ the path integral is taken between the cell values $\alpha_{1,i}$ and $\alpha_{1,i+1}$, also for second-order face
values, along the linear path with the face velocity frozen at the mean
$\bar u_{i+1/2}$, $\int_0^1u\,\partial_\ell\alpha_1\,\mathrm d\ell=\bar u_{i+1/2}\,(\alpha_{1,i+1}-\alpha_{1,i})$, and the fluctuation is split

% ===== main.tex, old lines 700-700, new lines 724-725
% OLD: the phase content. The shared factor is therefore limited further, for every phase and both functionals,
the phase content. The shared factor is therefore limited further, for every stiffened or Noble--Abel stiffened phase with $L_i>0$ and both
functionals,

% ===== main.tex, old lines 708-708, new lines 733-733
% OLD: The time step is given by the CFL condition of the advective sub-system,
The time step is given by the CFL condition of the explicit transport,

% ===== main.tex, old lines 721-721, new lines 746-747
% OLD: All contributions of the pressure--relaxation sub-system are treated implicitly and solved together, within a single step, as follows:
The pressure flux, the pressure work and the relaxation terms, together with the increments of the mass and internal energy
fluxes, are treated implicitly and solved together, within a single step, as follows:

% ===== main.tex, old lines 734-734, new lines 760-761
% OLD: With the incremental fluxes of Sec.~\ref{sec:split}, the semi-discrete form of \eqref{eq:pressure_sub} is
With the incremental fluxes of Sec.~\ref{sec:split}, the semi-discrete implicit step, \eqref{eq:pressure_sub} without the old-velocity part of
its internal energy fluxes and with the incremental mass fluxes, is

% ===== main.tex, old lines 810-810, new lines 837-838
% OLD: face pressure contributes $\tau\,\mathbb L(\eta^u,u)$, a bulk viscosity with coefficient $\tfrac12\rho c\Delta x$ for equal impedances; the
face pressure contributes $\tau\,\mathbb L(\eta^u,u^*)$, a bulk viscosity with coefficient $\tfrac12\rho c\Delta x$ for equal impedances, acting on
the velocity of $X^*$; the

% ===== main.tex, old lines 812-812, new lines 840-841
% OLD: \eqref{eq:wave} it contributes $\tau K^W\mathbb L(\eta^P,P)$, with $K^W\eta^P=\tfrac12c\Delta x$ for equal impedances and $c_p=c_f=c$. Both terms are implicit, so that they do not restrict the time step \eqref{eq:cfl}. In two dimensions the
\eqref{eq:wave} it contributes $\tau K^W\mathbb L(\eta^P,P)$, with $K^W\eta^P=\tfrac12c\Delta x$ for equal impedances and $c_p=c_f=c$. The pressure term is implicit; the velocity term is explicit,
with the diffusion number $\tfrac12\varphi_f\psi^u_fc\tau/\Delta x\le\tfrac18$ by \eqref{eq:weight}, so that neither restricts the time step \eqref{eq:cfl}. In two dimensions the

% ===== main.tex, old lines 825-825, new lines 854-854
% OLD: The face sweeps $u_f\tau$ during the step, of which the explicit flux has already swept $\bar u^e_f\tau$, $\bar u^e_f=\bar u_{i+1/2}$ of \eqref{eq:pc_rusanov}; the
The face sweeps $u_f\tau$ during the step, of which the explicit flux has already swept $\bar u^e_f\tau$, $\bar u^e_f=\bar u_{i+1/2}$ of \eqref{eq:rusanov_upwind}; the

% ===== main.tex, old lines 832-832, new lines 861-864
% OLD: with $(\cdot)^+=\max(\cdot,0)$ and $(\cdot)^-=\min(\cdot,0)$.
with $(\cdot)^-=\min(\cdot,0)$. Added to the explicit flux \eqref{eq:rusanov_upwind}, the increment gives the flux of the step,
$X^L_f(u_f)^++X^R_f(u_f)^--d_f(X^R_f-X^L_f)$ with $d_f$ the dissipation coefficient of the explicit flux given after \eqref{eq:rusanov_upwind}: $m_k$ and $\varepsilon_k$ are transported by the
face velocity $u_f$ from its upwind side with the face values of the explicit state, $\varepsilon_ku_f$ being the internal energy flux of the
pressure--relaxation sub-system, and the explicit step contributes their common dissipation.

% ===== main.tex, old lines 1030-1030, new lines 1062-1063
% OLD: jumps where $\delta(P)$ crosses a break point, by an amount that vanishes at a solution.
jumps where $\delta(P)$ crosses a break point, by an amount that vanishes at a solution of the closures with a common pressure
and is of order $sR(p_1-p_2)$ for finite $\mu$.

% ===== main.tex, old lines 1070-1071, new lines 1103-1106
% OLD: with $A_k=\varepsilon^{\rm tr}_k+\lambda_kD_0+\xi_k$, unique where $c_k^2>0$ on the bracket and found by a Newton iteration in $\alpha$ safeguarded by
% OLD: bisection in $\ln\alpha$. A stiffened or Noble--Abel stiffened partner phase is admitted only where $N_k>0$, and the outer
with $A_k=\varepsilon^{\rm tr}_k+\lambda_kD_0+\xi_k$, unique where $c_k^2>0$ on the bracket. The bracket is searched on the branch
$c_k^2>0$, on which $f_k$ is increasing: an initial volume with $c_k^2\le0$, as for a trace phase that has received mass and energy
without volume, is first doubled until $c_k^2>0$. The root is found by a Newton iteration in $\alpha$ safeguarded by bisection
in $\ln\alpha$. A stiffened or Noble--Abel stiffened partner phase is admitted only where $N_k>0$, and the outer

% ===== main.tex, old lines 1184-1184, new lines 1219-1219
% OLD: phases are compressed at constant specific internal energy, as in the advective sub-system \eqref{eq:advective}.
phases are compressed at constant specific internal energy, $m_k$ and $\varepsilon_k$ being transported by the same linear Rusanov operator with a shared limiter.

% ===== main.tex, old lines 1266-1266, new lines 1301-1301
% OLD: time linearisation introduced by the splitting, explicit velocity in the advective fluxes and implicit velocity
time linearisation introduced by the splitting, explicit velocity in the transport fluxes and implicit velocity

% ===== main.tex, old lines 1272-1272, new lines 1307-1307
% OLD: where $\mathcal H$ is the spatial discretisation of the split system, $\mathbf U_E$ the state on which the advective sub-system
where $\mathcal H$ is the spatial discretisation of the split system, $\mathbf U_E$ the state on which the explicit transport

% ===== main.tex, old lines 1291-1291, new lines 1326-1326
% OLD: $\mathbf U_{E,i}$, the update starts from $\tilde{\mathbf U}_{I,i}$ instead of $\mathbf U^n$, and the implicit step is $\tau=a_{ii}\Delta t$, giving the stage
$\mathbf U_{E,i}$, the update starts from $\tilde{\mathbf U}_{I,i}$ instead of $\mathbf U^n$, and the step, explicit transport included, has the length $\tau=a_{ii}\Delta t$, giving the stage

% ===== main.tex, old lines 1302-1302, new lines 1337-1337
% OLD: \begin{figure}[H]
\begin{figure}[p]

% ===== main.tex, old lines 1304-1307, new lines 1339-1342
% OLD: \resizebox*{!}{0.9\textheight}{% Flowchart of the semi-implicit scheme (Sec. 3). Needs: tikz with the libraries positioning, arrows.meta, calc, fit, backgrounds.
\begin{tikzpicture}[
  font=\fontsize{8pt}{9.6pt}\selectfont, >={Stealth[length=2.8mm,width=2.2mm]}, node distance=4mm,
  hdr/.style={fill=hdrblue, text=white, font=\fontsize{8pt}{9.6pt}\selectfont\bfseries, inner sep=2.5pt, rounded corners=1pt,
              align=flush center},
  eq/.style={draw=black!40, fill=white, inner sep=3pt, align=flush center},
  outer/.style={draw=hdrblue!80, fill=hdrblue!8, rounded corners=4pt, inner sep=5pt},
  step/.style={draw=black!45, fill=black!5, align=flush center, inner sep=3pt, text width=#1},
  step/.default=9.6cm,
  note/.style={align=flush center, font=\fontsize{7pt}{8.4pt}\selectfont},
  dashedarrow/.style={->, dashed, semithick, black!75},
]
\definecolor{hdrblue}{RGB}{31,73,125}
\definecolor{clgreen}{RGB}{214,232,220}
\definecolor{rlpurple}{RGB}{226,222,240}

% ----------------------------------------------------------------- full system
\node[eq, text width=9.6cm] (full) {$\partial_t\mathbf U+\partial_x\mathbf F^c(\mathbf U)+\mathbf B^c(\mathbf U)\,\partial_x\mathbf U
  +\partial_x\mathbf F^p(\mathbf U)+\mathbf B^p(\mathbf U)\,\partial_xu=\mathbf S(\mathbf U)$};
\node[hdr, above=0pt of full.north, text width=9.6cm] (fullh) {Full system};
\begin{scope}[on background layer]
  \node[outer, fit=(full)(fullh)] (fullbox) {};
\end{scope}

\node[hdr, below=6mm of fullbox] (split) {Flux Vector Splitting};
\draw[thick, black!65] (fullbox.south) -- (split.north);

% ----------------------------------------------------------------- the two sub-systems
\node[eq, below left=15mm and 2.5mm of split.south, anchor=north east, text width=4.75cm, minimum height=1.9cm] (adv)
  {$\partial_t\mathbf U+\partial_x\mathbf F^c+\mathbf B^c\,\partial_x\mathbf U=\mathbf 0$\\[3pt]
   {\fontsize{7pt}{8.4pt}\selectfont path-conservative Rusanov,\\ material speed, flux $\mathbf F^u$;\\ $\alpha_1$ advected entirely}};
\node[hdr, above=0pt of adv.north, text width=4.75cm] (advh) {Advective, explicit};
\node[eq, below right=15mm and 2.5mm of split.south, anchor=north west, text width=4.75cm, minimum height=1.9cm] (prs)
  {$\partial_t\mathbf U+\partial_x\mathbf F^p+\mathbf B^p\,\partial_xu=\mathbf S(\mathbf U)$\\[3pt]
   {\fontsize{7pt}{8.4pt}\selectfont $\varepsilon_ku$, $m_ku$ as increments $X(u^{n+1}-u^n)$;\\ pressure work $+$ relaxations\\ in one local closure}};
\node[hdr, above=0pt of prs.north, text width=4.75cm] (prsh) {Pressure--relaxation, implicit};
\begin{scope}[on background layer]
  \node[outer, fit=(adv)(advh)] (advbox) {};
  \node[outer, fit=(prs)(prsh)] (prsbox) {};
\end{scope}
\draw[thick, black!65] (split.south west) -- (advbox.north);
\draw[thick, black!65] (split.south east) -- (prsbox.north);

% ----------------------------------------------------------------- one stage (steps numbered as in Sec. 3.6)
\node[hdr, below=11mm of split.south |- advbox.south, text width=9.6cm] (alh) {Algorithm: one stage};
\node[step, below=4mm of alh] (s0)
  {\textbf{Before the first stage of a step}\\ activation of the closures; projection of $\mathbf U^n$ (instantaneous $pT$, $pTg$)};
\node[step, below=5mm of s0] (s1)
  {\textbf{1. Explicit transport} on $\mathbf U^e$, from $\mathbf U^{\rm b}$\\
   $X^*$: Rusanov fluxes, \ $\alpha^*_1=\alpha^{\rm b}_1-r_x\big(\delta_x[F^e[\alpha_1]]-\alpha^e_1\,\delta_x[\bar u^e]\big)$};
\node[step, below=5mm of s1] (s2)
  {\textbf{2. Face coefficients} on $X^*$, frozen over the stage\\ $u_f$, $p_f$ affine in the adjacent $P$; weight $\varphi$, factors $\psi$};
\node[step, below=5mm of s2] (s3)
  {\textbf{Newton iteration for the cell pressures $P$}\\[1pt]
   3. face velocity with implicit pressure Laplace term;\\
   \phantom{3. }face pressure with explicit velocity Laplace term on $u^*$;\\
   \phantom{3. }upwind incremental fluxes $\Rightarrow$ $X^{\rm tr}$, $s$, $\rho\mathbf u$, $E$, $(\rho e)^C$\\
   4. local closure $\Rightarrow$ $\bar p^{\rm cl}$, phase states\\
   5. Newton update of $G(P)=\bar p^{\rm cl}-P$\\[1pt]
   \textbf{6. Update:} $\rho\mathbf u$, $E$ from 3; $\alpha_k$, $m_k$, $\varepsilon_k$ from 4};
\begin{scope}[on background layer]
  \node[outer, fit=(alh)(s0)(s1)(s2)(s3)] (algbox) {};
\end{scope}
\draw[thick, black!65] (advbox.south) -- (advbox.south |- algbox.north);
\draw[thick, black!65] (prsbox.south) -- (prsbox.south |- algbox.north);
\draw[thick, black!65] (alh) -- (s0); \draw[thick, black!65] (s0) -- (s1); \draw[thick, black!65] (s1) -- (s2); \draw[thick, black!65] (s2) -- (s3);

% ----------------------------------------------------------------- local closure with energy partition
\node[step=6.3cm, fill=clgreen, draw=black!45, below left=14mm and 6mm of algbox.south, anchor=north east, minimum height=4.3cm] (cl)
  {\textbf{4. Local closure} at every residual evaluation\\[2pt]
   \textbf{energy defect}\\ $\delta=(\rho e)^C-\sum_k\varepsilon^{\rm tr}_k+s\sum_k\alpha^*_kp_k$\\[2pt]
   \textbf{partition} $\Lambda_k(\delta)$, $\sum_k\Lambda_k=\delta$:\\
   $0\le\delta\le H$: shock share $H_k/H$, $H_k$ on $X^*$\\
   excess over $H$ and $\delta<0$: $\lambda^{\mathcal E}_k=\mathcal E_k/\sum_j\mathcal E_j$\\[2pt]
   $\Rightarrow\ \sum_k\varepsilon_k=(\rho e)^C$ exactly};

% ----------------------------------------------------------------- relaxation closures by rates
\node[step=6.3cm, fill=rlpurple, draw=black!45, below right=14mm and 6mm of algbox.south, anchor=north west, minimum height=4.3cm] (rl)
  {\textbf{Relaxation closure by rates}\\
   ordering: $\mu=\infty$ if $\theta>0$ or $\nu>0$\\ \phantom{ordering:} $\theta=\infty$ if $\nu>0$\\[3pt]
   one formula per process, $\zeta=1/(\tau\,{\rm rate})\in[0,\infty]$:\\[2pt]
   \begin{tabular}{@{}l@{\ \ }l@{}}
   $\nu>0$:      & $\zeta\,(\rho Y-m^{\rm tr}_2)+g_2-g_1=0$\\[2pt]
   $\theta>0$:   & $p$ closure $+$ heat $\tau\mathcal Q=\dfrac{T^0_2-T^0_1}{\sum_kC_{pk}^{-1}+\zeta}$\\[6pt]
   otherwise:    & $\zeta R=p_1-p_2$
   \end{tabular}\\[3pt]
   $\zeta=0$: $pTg$, $pT$, $p$ closure;\ \ $\zeta=\infty$: off\\[2pt]
   $\Rightarrow$ one scalar equation per cell:\\ $\Phi(p)=0$, or $\mathcal F(R)=0$ for finite $\mu$};

\draw[dashedarrow] (s3.south -| cl.north) -- node[pos=0.55, left, note] {$s$, $X^{\rm tr}$,\\$(\rho e)^C$} (cl.north);
\draw[dashedarrow] (cl.east) -- node[above, note] {$\mu,\theta,\nu$} (rl.west);
\draw[dashedarrow] (rl.north) -- node[pos=0.45, right, note] {$\bar p^{\rm cl}$,\\phase states} (s3.south -| rl.north);

\end{tikzpicture}
}
% OLD: \caption{Schematic of the semi-implicit scheme. The Zha--Bilgen splitting separates the advective
% OLD: sub-system, advanced explicitly by the path-conservative Rusanov scheme with the material speed, from the
% OLD: pressure--relaxation sub-system, solved implicitly together with the incremental fluxes. Every residual evaluation of the pressure system calls the local closure of each cell, which
\resizebox*{!}{0.88\textheight}{% Flowchart of the semi-implicit scheme (Sec. 3). Needs: tikz with the libraries positioning, arrows.meta, calc, fit, backgrounds.
\begin{tikzpicture}[
  font=\fontsize{8pt}{9.6pt}\selectfont, >={Stealth[length=2.8mm,width=2.2mm]}, node distance=4mm,
  hdr/.style={fill=hdrblue, text=white, font=\fontsize{8pt}{9.6pt}\selectfont\bfseries, inner sep=2.5pt, rounded corners=1pt,
              align=flush center},
  eq/.style={draw=black!40, fill=white, inner sep=3pt, align=flush center},
  outer/.style={draw=hdrblue!80, fill=hdrblue!8, rounded corners=4pt, inner sep=5pt},
  step/.style={draw=black!45, fill=black!5, align=flush center, inner sep=3pt, text width=#1},
  step/.default=9.6cm,
  note/.style={align=flush center, font=\fontsize{7pt}{8.4pt}\selectfont},
  dashedarrow/.style={->, dashed, semithick, black!75},
]
\definecolor{hdrblue}{RGB}{31,73,125}
\definecolor{clgreen}{RGB}{214,232,220}
\definecolor{rlpurple}{RGB}{226,222,240}

% ----------------------------------------------------------------- full system
\node[eq, text width=9.6cm] (full) {$\partial_t\mathbf U+\partial_x\mathbf F^c(\mathbf U)+\mathbf B^c(\mathbf U)\,\partial_x\mathbf U
  +\partial_x\mathbf F^p(\mathbf U)+\mathbf B^p(\mathbf U)\,\partial_xu=\mathbf S(\mathbf U)$};
\node[hdr, above=0pt of full.north, text width=9.6cm] (fullh) {Full system};
\begin{scope}[on background layer]
  \node[outer, fit=(full)(fullh)] (fullbox) {};
\end{scope}

\node[hdr, below=6mm of fullbox] (split) {Flux Vector Splitting};
\draw[thick, black!65] (fullbox.south) -- (split.north);

% ----------------------------------------------------------------- the two sub-systems
\node[eq, below left=15mm and 2.5mm of split.south, anchor=north east, text width=4.75cm, minimum height=1.9cm] (adv)
  {$\partial_t\mathbf U+\partial_x\mathbf F^c+\mathbf B^c\,\partial_x\mathbf U=\mathbf 0$\\[3pt]
   {\fontsize{7pt}{8.4pt}\selectfont path-conservative Rusanov,\\ material speed, flux $\mathbf F^u$;\\ $\alpha_1$ advected entirely}};
\node[hdr, above=0pt of adv.north, text width=4.75cm] (advh) {Advective, explicit};
\node[eq, below right=15mm and 2.5mm of split.south, anchor=north west, text width=4.75cm, minimum height=1.9cm] (prs)
  {$\partial_t\mathbf U+\partial_x\mathbf F^p+\mathbf B^p\,\partial_xu=\mathbf S(\mathbf U)$\\[3pt]
   {\fontsize{7pt}{8.4pt}\selectfont $\varepsilon_ku$, $m_ku$ as increments $X(u^{n+1}-u^n)$;\\ pressure work $+$ relaxations\\ in one local closure}};
\node[hdr, above=0pt of prs.north, text width=4.75cm] (prsh) {Pressure--relaxation, implicit};
\begin{scope}[on background layer]
  \node[outer, fit=(adv)(advh)] (advbox) {};
  \node[outer, fit=(prs)(prsh)] (prsbox) {};
\end{scope}
\draw[thick, black!65] (split.south west) -- (advbox.north);
\draw[thick, black!65] (split.south east) -- (prsbox.north);

% ----------------------------------------------------------------- one stage (steps numbered as in Sec. 3.6)
\node[hdr, below=11mm of split.south |- advbox.south, text width=9.6cm] (alh) {Algorithm: one stage};
\node[step, below=4mm of alh] (s0)
  {\textbf{Before the first stage of a step}\\ activation of the closures; projection of $\mathbf U^n$ (instantaneous $pT$, $pTg$)};
\node[step, below=5mm of s0] (s1)
  {\textbf{1. Explicit transport} on $\mathbf U^e$, from $\mathbf U^{\rm b}$\\
   $X^*$: Rusanov fluxes, \ $\alpha^*_1=\alpha^{\rm b}_1-r_x\big(\delta_x[F^e[\alpha_1]]-\alpha^e_1\,\delta_x[\bar u^e]\big)$};
\node[step, below=5mm of s1] (s2)
  {\textbf{2. Face coefficients} on $X^*$, frozen over the stage\\ $u_f$, $p_f$ affine in the adjacent $P$; weight $\varphi$, factors $\psi$};
\node[step, below=5mm of s2] (s3)
  {\textbf{Newton iteration for the cell pressures $P$}\\[1pt]
   3. face velocity with implicit pressure Laplace term;\\
   \phantom{3. }face pressure with explicit velocity Laplace term on $u^*$;\\
   \phantom{3. }upwind incremental fluxes $\Rightarrow$ $X^{\rm tr}$, $s$, $\rho\mathbf u$, $E$, $(\rho e)^C$\\
   4. local closure $\Rightarrow$ $\bar p^{\rm cl}$, phase states\\
   5. Newton update of $G(P)=\bar p^{\rm cl}-P$\\[1pt]
   \textbf{6. Update:} $\rho\mathbf u$, $E$ from 3; $\alpha_k$, $m_k$, $\varepsilon_k$ from 4};
\begin{scope}[on background layer]
  \node[outer, fit=(alh)(s0)(s1)(s2)(s3)] (algbox) {};
\end{scope}
\draw[thick, black!65] (advbox.south) -- (advbox.south |- algbox.north);
\draw[thick, black!65] (prsbox.south) -- (prsbox.south |- algbox.north);
\draw[thick, black!65] (alh) -- (s0); \draw[thick, black!65] (s0) -- (s1); \draw[thick, black!65] (s1) -- (s2); \draw[thick, black!65] (s2) -- (s3);

% ----------------------------------------------------------------- local closure with energy partition
\node[step=6.3cm, fill=clgreen, draw=black!45, below left=14mm and 6mm of algbox.south, anchor=north east, minimum height=4.3cm] (cl)
  {\textbf{4. Local closure} at every residual evaluation\\[2pt]
   \textbf{energy defect}\\ $\delta=(\rho e)^C-\sum_k\varepsilon^{\rm tr}_k+s\sum_k\alpha^*_kp_k$\\[2pt]
   \textbf{partition} $\Lambda_k(\delta)$, $\sum_k\Lambda_k=\delta$:\\
   $0\le\delta\le H$: shock share $H_k/H$, $H_k$ on $X^*$\\
   excess over $H$ and $\delta<0$: $\lambda^{\mathcal E}_k=\mathcal E_k/\sum_j\mathcal E_j$\\[2pt]
   $\Rightarrow\ \sum_k\varepsilon_k=(\rho e)^C$ exactly};

% ----------------------------------------------------------------- relaxation closures by rates
\node[step=6.3cm, fill=rlpurple, draw=black!45, below right=14mm and 6mm of algbox.south, anchor=north west, minimum height=4.3cm] (rl)
  {\textbf{Relaxation closure by rates}\\
   ordering: $\mu=\infty$ if $\theta>0$ or $\nu>0$\\ \phantom{ordering:} $\theta=\infty$ if $\nu>0$\\[3pt]
   one formula per process, $\zeta=1/(\tau\,{\rm rate})\in[0,\infty]$:\\[2pt]
   \begin{tabular}{@{}l@{\ \ }l@{}}
   $\nu>0$:      & $\zeta\,(\rho Y-m^{\rm tr}_2)+g_2-g_1=0$\\[2pt]
   $\theta>0$:   & $p$ closure $+$ heat $\tau\mathcal Q=\dfrac{T^0_2-T^0_1}{\sum_kC_{pk}^{-1}+\zeta}$\\[6pt]
   otherwise:    & $\zeta R=p_1-p_2$
   \end{tabular}\\[3pt]
   $\zeta=0$: $pTg$, $pT$, $p$ closure;\ \ $\zeta=\infty$: off\\[2pt]
   $\Rightarrow$ one scalar equation per cell:\\ $\Phi(p)=0$, or $\mathcal F(R)=0$ for finite $\mu$};

\draw[dashedarrow] (s3.south -| cl.north) -- node[pos=0.55, left, note] {$s$, $X^{\rm tr}$,\\$(\rho e)^C$} (cl.north);
\draw[dashedarrow] (cl.east) -- node[above, note] {$\mu,\theta,\nu$} (rl.west);
\draw[dashedarrow] (rl.north) -- node[pos=0.45, right, note] {$\bar p^{\rm cl}$,\\phase states} (s3.south -| rl.north);

\end{tikzpicture}
}
\caption{Schematic of the semi-implicit scheme. The Toro--V\'azquez-Cend\'on splitting separates the advective
sub-system, advanced explicitly with the old-velocity part of the internal energy fluxes, from the pressure--relaxation
sub-system, solved implicitly with the increments of the mass and internal energy fluxes. Every residual evaluation of the pressure system calls the local closure of each cell, which

% ===== main.tex, old lines 1322-1322, new lines 1357-1357
% OLD: \item \emph{Explicit advective sub-system.} The face values are reconstructed on $\mathbf U^e$ by \eqref{eq:shared_limiter} with the cap \eqref{eq:phi_cap}, and every
\item \emph{Explicit transport.} The face values are reconstructed on $\mathbf U^e$ by \eqref{eq:shared_limiter} with the cap \eqref{eq:phi_cap}, and every

% ===== main.tex, old lines 1343-1343, new lines 1378-1378
% OLD: with $\Theta_f=\tau/(\bar\rho_f h)+\eta^P_f/h$, $W_f=\rho^*_R/(\rho^*_L+\rho^*_R)$ and $\varpi_f=\eta^u_fD^*_f/h$, $h=\Delta x$ or $\Delta y$. At a boundary the
with $\Theta_f=\tau/(\bar\rho_f h)+\eta^P_f/h$, $W_f=\rho^*_R/(\rho^*_L+\rho^*_R)$ and $\varpi_f=\eta^u_fD^*_f/h$, $h=\Delta x$ or $\Delta y$. At a non-periodic boundary the

% ===== main.tex, old lines 1345-1345, new lines 1380-1380
% OLD: boundaries and their mirror image at walls, with the normal velocity $2u_w-u$ for a wall velocity $u_w$; on a wall face $u_f=u_w$.
boundaries and their mirror image at walls, with the normal velocity $2u_w-u$ for a wall velocity $u_w$; on a wall face $u_f=u_w$. At periodic boundaries the ghost cells and their pressures are those of the opposite side.

% ===== main.tex, old lines 1408-1408, new lines 1443-1443
% OLD: Sec.~\ref{sec:pressure_system}; the first time step is an acoustic step with Courant number $0.4$. Thermal and chemical relaxation
Sec.~\ref{sec:pressure_system} and the start-up rule of Sec.~\ref{sec:time}. Thermal and chemical relaxation

% ===== main.tex, old lines 1459-1459, new lines 1494-1494
% OLD: otherwise, $t_{\rm end}=0.25$\,s ($M=|u|/c$ between $2.7\times10^{-4}$ in the air and $6\times10^{-4}$ in the water); and a single gas, air, with a density step $\rho=1000$ for $x<0.25$
otherwise, $t_{\rm end}=0.25$\,s, so that $M=|u|/c$ lies between $2.7\times10^{-4}$ in the air and $6\times10^{-4}$ in the water; and a single gas, air, with a density step $\rho=1000$ for $x<0.25$

% ===== main.tex, old lines 1492-1492, new lines 1527-1527
% OLD: The following Riemann problems on $[0,1]$ (on $[-10,10]$ for the detonation products) with transmissive boundaries
The following Riemann problems on $[0,1]$, or $[-10,10]$ for the detonation products, with transmissive boundaries

% ===== main.tex, old lines 1533-1533, new lines 1568-1568
% OLD: Dodecane is NASG with liquid $(1.206,\ 1.681\times10^8,\ -996054,\ 0,\ 2532,\ 7.51\times10^{-4})$ and vapour
Dodecane is NASG with liquid $(1.206,\allowbreak\ 1.681\times10^8,\allowbreak\ -996054,\allowbreak\ 0,\allowbreak\ 2532,\allowbreak\ 7.51\times10^{-4})$ and vapour

% ===== main.tex, old lines 1566-1567, new lines 1601-1602
% OLD: a stationary solution of the incompressible Euler equations \cite{TaylorGreen1937}. The volume fraction is $\alpha_2=\tfrac12$ (single-phase
% OLD: case) or $\alpha_2=\tfrac12+\tfrac1{10}(\cos2x+\cos2y)$ (two-phase case). The background pressure
a stationary solution of the incompressible Euler equations \cite{TaylorGreen1937}. The volume fraction is $\alpha_2=\tfrac12$ in the single-phase
case and $\alpha_2=\tfrac12+\tfrac1{10}(\cos2x+\cos2y)$ in the two-phase case. The background pressure

% ===== main.tex, old lines 1589-1589, new lines 1624-1624
% OLD: bottom and top boundaries are walls and the right boundary is transmissive. $300\times100$ cells (the grid of Pelanti and Shyue \cite[Sec.~5.3]{PelantiShyue2014}), $t_{\rm end}=8.6$\,ms,
bottom and top boundaries are walls and the right boundary is transmissive. $300\times100$ cells, the grid of Pelanti and Shyue \cite[Sec.~5.3]{PelantiShyue2014}, $t_{\rm end}=8.6$\,ms,

% ===== main.tex, old lines 1598-1599, new lines 1633-1634
% OLD: $\alpha_v=10^{-6}$ outside, at rest. All boundaries are walls; the left wall moves with velocity $100$\,m/s in $+x$ (the grid
% OLD: is fixed, and the wall boundary condition imposes this normal velocity). $200\times200$ cells, $t_{\rm end}=1.2$\,ms,
$\alpha_v=10^{-6}$ outside, at rest. All boundaries are walls; the left wall moves with velocity $100$\,m/s in $+x$; the grid
is fixed and the wall boundary condition imposes this normal velocity. $200\times200$ cells, $t_{\rm end}=1.2$\,ms,

% ===== main.tex, old lines 1609-1609, new lines 1644-1644
% OLD: boundaries elsewhere. $361\times200$ cells (half the resolution of \cite{Pelanti2022}), $t_{\rm end}=0.4$\,ms, $p$ relaxation; the pressure is recorded at the wall
boundaries elsewhere. $361\times200$ cells, half the resolution of \cite{Pelanti2022}, $t_{\rm end}=0.4$\,ms, $p$ relaxation; the pressure is recorded at the wall

% ===== main.tex, old lines 1619-1620, new lines 1654-1655
% OLD: two-phase model with total energy and mechanical, thermal and chemical relaxation, for stiffened-gas,
% OLD: Noble--Abel stiffened-gas and Mie--Gr\"uneisen equations of state, in one and two space dimensions.
two-phase model with total energy, for stiffened-gas and Noble--Abel stiffened-gas phases with mechanical, thermal and
chemical relaxation and for Mie--Gr\"uneisen phases with instantaneous mechanical relaxation, in one and two space dimensions.

% ===== main.tex, old lines 1624-1628, new lines 1659-1663
% OLD: finite or infinite rate, a slower process being active only when the faster ones are instantaneous \cite{Pelanti2022}, enter one local closure per cell, which reduces to a scalar equation for the
% OLD: cell pressure; each stage then reduces to a mildly nonlinear system for the cell pressures with a three-point
% OLD: (five-point in 2D) Jacobian, solved by Newton's method. The time step is restricted by the flow velocity only, whatever
% OLD: the Mach number and the relaxation rates. We have proved that, under non-degeneracy assumptions, the scheme converges
% OLD: to the corresponding equilibrium scheme as each relaxation rate tends to infinity and reduces to the scheme without that
finite or infinite rate, a slower process being active only when the faster ones are instantaneous \cite{Pelanti2022}, enter one local closure per cell, which reduces to a scalar equation in the
cell pressure or, for finite mechanical relaxation, in the volume increment; each stage then reduces to a nonlinear system
for the cell pressures with a three-point Jacobian, five-point in 2D, solved by Newton's method. The time step is restricted by the flow velocity only, whatever
the Mach number and the relaxation rates. We have proved that, at fixed mesh and under non-degeneracy assumptions, the scheme converges, for data on the
equilibrium manifold, to the corresponding equilibrium scheme as each relaxation rate tends to infinity and reduces to the scheme without that

% ===== main.tex, old lines 1634-1635, new lines 1669-1670
% OLD: reduces to the unique consistent local share, and the discretisation error with the unique share that keeps both phases
% OLD: admissible for every defect. No calibrated heat exchange is needed.
reduces to the unique consistent local share, and the discretisation error with the unique linear share that keeps both phases
admissible whenever possible. No calibrated heat exchange is needed.

% ===== main.tex, old lines 1637-1642, new lines 1672-1679
% OLD: The flux is split following Zha and Bilgen: all quantities are transported explicitly with the old velocity by the
% OLD: path-conservative Rusanov scheme, whose dissipation speed is the material speed, the volume fraction being advected
% OLD: entirely in this step, while the pressure, its work on each phase and an incremental flux of the conservative transported
% OLD: quantities are treated implicitly. The central implicit operator is supplemented by implicit Laplace dissipation of the
% OLD: pressure and the velocity with acoustic coefficients, weighted by the deficit of the implicit viscosity with respect to
% OLD: the upwind one, so that it vanishes in the low Mach number regime, and limited to first order at local extrema while second order in smooth regions. Second
The flux is split following Toro and V\'azquez-Cend\'on, with the mass and internal energy fluxes evaluated at the new
velocity: all quantities are transported explicitly with the old velocity by the path-conservative Rusanov scheme, whose
dissipation speed is the material speed, the volume fraction being advected entirely in this step, while the pressure,
its work on each phase and the increments of the mass and internal energy fluxes from the old to the new velocity are
treated implicitly. The central implicit operator is supplemented by Laplace dissipation of the
pressure and the velocity with acoustic coefficients, weighted by the deficit of the backward-Euler viscosity with respect to
the upwind one, so that it vanishes in the low Mach number regime; it is first order at local extrema and second order in
smooth regions. Second

% ===== main.tex, old lines 1644-1644, new lines 1681-1682
% OLD: preserves moving contacts and the phase temperatures under transport.
preserves moving contacts between stiffened-gas or Noble--Abel stiffened-gas phases, at uniform temperature when thermal or
chemical relaxation is active.

% ===== main.tex, old lines 1646-1648, new lines 1684-1687
% OLD: The numerical results confirm these properties: second-order convergence on a smooth two-phase vortex, a Newton
% OLD: iteration count and a time step independent of the Mach number, face velocities that are discretely divergence-free to
% OLD: $O(M^2)$ down to $M\approx10^{-6}$ on the Taylor--Green vortex, and agreement with exact solutions and reference results
The numerical results confirm these properties: second-order convergence of the velocity and the volume fraction on a
smooth two-phase vortex at Mach numbers from $2.4$ to $0.024$, and of the pressure for well-prepared data, face velocities
that are discretely divergence-free to $O(M^2)$ down to $M\approx10^{-6}$ on the Taylor--Green vortex with a time step
independent of the Mach number, and agreement with exact solutions and reference results

% ===== main.tex, old lines 1650-1650, new lines 1689-1689
% OLD: two-dimensional shock--bubble, Richtmyer--Meshkov, bubble-collapse and underwater-explosion problems, at instantaneous
two-dimensional shock--bubble, Richtmyer--Meshkov, vapour-bubble compression and underwater-explosion problems, at instantaneous

% ===== figs/scheme.tex, old lines 31-31, new lines 31-31
% OLD:    {\fontsize{7pt}{8.4pt}\selectfont path-conservative Rusanov,\\ material speed;\\ $\alpha_1$ advected entirely}};
   {\fontsize{7pt}{8.4pt}\selectfont path-conservative Rusanov,\\ material speed, flux $\mathbf F^u$;\\ $\alpha_1$ advected entirely}};

% ===== figs/scheme.tex, old lines 35-35, new lines 35-35
% OLD:    {\fontsize{7pt}{8.4pt}\selectfont $+$ incremental fluxes $X(u^{n+1}-u^n)$;\\ pressure work $+$ relaxations\\ in one local closure}};
   {\fontsize{7pt}{8.4pt}\selectfont $\varepsilon_ku$, $m_ku$ as increments $X(u^{n+1}-u^n)$;\\ pressure work $+$ relaxations\\ in one local closure}};

% ===== figs/scheme.tex, old lines 49-49, new lines 49-49
% OLD:   {\textbf{1. Explicit advective sub-system} on $\mathbf U^e$, from $\mathbf U^{\rm b}$\\
  {\textbf{1. Explicit transport} on $\mathbf U^e$, from $\mathbf U^{\rm b}$\\

% ===== figs/scheme.tex, old lines 55-55, new lines 55-56
% OLD:    3. face velocity and pressure with implicit Laplace dissipation;\\
   3. face velocity with implicit pressure Laplace term;\\
   \phantom{3. }face pressure with explicit velocity Laplace term on $u^*$;\\

